\section{本章纲要：旧秩序的语言，新力量的竞争}

东周王室保住了“天下”的名分，诸侯则掌握越来越多的实际资源。齐、晋、楚的霸权证明，会盟和礼仪能够暂时组织秩序；晋的卿族、楚的扩张、吴越的进入又证明，这种秩序无法冻结国家内部和边缘地区的变化。战国将把这些变化推进到更直接的制度竞争中。

\subsection{把春秋读成一组问题，而不是一串霸主}

本章的第一条线，是王室名分与可调动资源的分离。\eventref{SA-0770-EAST}{春秋·东迁}之后，周仍是诸侯可以援引的中心：它能够使继承、会盟和出兵拥有一种大家认识的语言；但它已很难单独确保道路、粮食和军队会按这一语言行动。\eventref{SA-0707-XUGE}{春秋·繻葛之战}将这种处境公开化。王师的败绩不是周礼从此无人理会，而是提醒人们，名分若不经过诸侯、卿族、城邑和供给网络，便不能自动成为力量。此后齐、晋、楚反复争取的，正是在危机中把资源与名分重新接合的资格。

第二条线，是联盟既能保护小国，也会使小国承担无法独自决定的成本。齐在\eventref{SA-0659-QI-XING}{春秋·齐迁邢于夷仪}中的救援、在\eventref{SA-0656-SHAOLING}{春秋·召陵之盟}中的北上，晋在\eventref{SA-0632-CHENGPU}{春秋·城濮之战}后的组织，楚在\eventref{SA-0597-BI}{春秋·邲之战}后的北方压力，都说明会盟不是纯粹仪式。它可能提供道路、援军与协商的场所，也会要求参盟者交出粮秣、时间和未来行动的余地。盟会没有把各国变成同一国家的郡县，正因如此，霸主必须不断以新的救援、惩戒、礼序或让步证明自己仍值得被跟随。

第三条线，是国家的能力为什么会同时造成国家内部的困难。晋能把军政职位、封邑和外交关系交给卿族，因而拥有持续调动诸侯的条件；这些位置也会在家门、家臣与地方城邑中积累，终于使\eventref{SA-0497-JIN-FACTION}{春秋·范氏中行氏赵氏之争}一类冲突不再只是朝廷内部的口角。齐的\eventref{SA-0548-CUIZHU}{春秋·崔杼杀齐庄公}、鲁的\eventref{SA-0517-LU-ZHAO-EXILE}{春秋·鲁昭公出奔}和楚对贵族网络的长期依赖，同样说明国君并不能越过一切中介直接掌握人、地与兵。家门并不只是国家的敌人；它们原本就是征发、守备与治理得以完成的渠道。问题在于，渠道一旦能把资源留在自身，公室便要面对谁有权代表国家、谁应为公共后果负责的新难题。

第四条线，是地理使每一种政治选择都有不同的价格。齐的海岱腹地和东部小国关系，晋的高原、河谷与中原通道，楚的江汉网络，秦的关中山口，吴越的水网，都不是地图上的背景色。它们决定军队怎样越过边界、粮食怎样接续、迁徙者在哪里落脚，也决定一国为何能把胜利送到远方、却难以长久占有陌生腹地。\eventref{SA-0506-BOJU}{春秋·柏举之战}之后吴军入郢，证明下游国家能够突破旧有的竞争格局；\eventref{SA-0505-YUE-INVADE-WU}{春秋·越军入吴}又提示，远征越深，本土守备越容易成为另一个战场。地理提供条件，绝不替人物作出决定；但不看地理，人物的选择也会失去它们实际承受的限制。

最后，本章将文本的保存方式视为历史的一部分。经文的短记、\sourcebook{左传}的叙事、\sourcebook{国语}的国别说辞、\sourcebook{史记}的汉代整合以及盟书、铜器、城址等物质材料，各自给出不同的尺度。它们可以互相照亮，不能相互取代。\eventref{SA-0481-END}{春秋·记事终点}只是鲁国传世编年的边界，而非所有行动突然停顿的证据；\eventref{SA-0473-WU-FALL}{春秋·越灭吴}和\eventref{WQ-0453-ZHI}{战国·三家灭智}则把读者带向仍在继续的资源重组。这样读春秋，读到的不是一条注定通往统一帝国的直线，而是一批反复被提出、由战国各国以更直接的征发和官僚办法回答的问题。

\begin{turningpoint}[本章留下的不是答案，而是国家化的压力]
春秋的得失不宜被归结为“分封必败”或“霸主无用”。封邑、卿族、会盟和礼仪在当时使跨区域合作成为可能；它们也让权力必须通过很多可被继承、争夺和撤回的关系才能抵达地方。战国的改革会试图缩短这条链条，把土地、兵役、官职和法律更直接接到政治中心；它会获得新的动员速度，也会让家庭与聚落更直接承担国家竞争的代价。后面的章节应当从这一压力继续观察，而不把它误当作已经在春秋完成的制度事实。
\end{turningpoint}

\subsection{给读者的一条阅读路径：先看行动，再问制度}

阅读春秋时，可以先抓住一次行动的最小骨架，再逐层追问它为何发生、怎样被写下、留下什么没有解决的压力。比如从\eventref{SA-0656-SHAOLING}{春秋·召陵之盟}出发，先确认齐与诸侯北抵楚境、双方以盟收束；随后再问齐为何能让江、黄、蔡、陈等处进入这场行动，楚为何愿意以交涉而非继续作战回应，以及苞茅、昭王等说辞在《左传》中怎样把一场政治互动塑造成礼制冲突。这样读，不必先相信一个“齐桓公称霸”的整齐结论，也不会因为叙事细节难以独立核实而否认会盟确实改变了各国的选择空间。

同样，可以从\eventref{SA-0597-BI}{春秋·邲之战}与\eventref{SA-0546-MIBING}{春秋·弭兵之会}之间的距离理解晋楚竞争。邲使晋的保护承诺遭到严厉检验，弭兵则表明长期动员、过境与供给已令各方需要降低直接冲突。二者之间并不存在“晋衰、楚盛”一条足以解释一切的直线：国君继承、卿族位置、小国改从、道路远近和军粮承受力都在变化。所谓制度得失，正是在这些改变中显出：一项能降低眼前风险的安排，可能把新的资源积累机会交给原本就占有职位和封邑的人。

还可以把\eventref{SA-0506-BOJU}{春秋·柏举之战}同\eventref{SA-0473-WU-FALL}{春秋·越灭吴}连读。前者说明吴能把江淮水陆网络和蔡、唐等关系转成深入楚腹地的远征；后者说明一次远征的光荣不足以替代长期的本土防务、朝廷判断和持续供给。二者之间的夫椒、艾陵、黄池以及越的恢复不应被压缩为“卧薪尝胆”的道德寓言。读者每次遇到后出文本中格外精彩的谋略、预言或人物姿态，都可先问：这是经年所能确认的行动，还是后来叙事用以解释行动的语言？这个问题不会使故事失色，反而使故事中的选择、信息不足和风险承担更为清楚。

最后，应把本章所见的家门政治同战国的改革问题相连，而不把它们混为一事。\eventref{SA-0497-JIN-FACTION}{春秋·范氏中行氏赵氏之争}显示城邑、部众与卿族关系已能脱离共同指挥彼此对抗；\eventref{WQ-0453-ZHI}{战国·三家灭智}和\eventref{WQ-0403-THREEJIN}{战国·三家分晋}才在后来的时间里让这种分裂获得更清楚的结果与名分。战国变法不是对春秋无能的简单否定，而是试图以更直接的登记、官职、军功和征发来处理春秋留下的中介层次。它同样会产生新的代价，因此本书接下来的问题不是“国家是否变强”这一句，而是谁被纳入新的动员链条，谁能检查权力的使用，谁又在新秩序中失去原先可以依靠的位置。
